\documentclass[11pt]{article}
\usepackage{amsmath,amsthm,amssymb}
\usepackage[margin=1in]{geometry}
\usepackage{stmaryrd}
\usepackage{array}
\usepackage[colorlinks=true,linkcolor=blue,citecolor=blue,urlcolor=blue]{hyperref}

\newtheorem{theorem}{Theorem}
\newtheorem{proposition}[theorem]{Proposition}
\newtheorem{corollary}[theorem]{Corollary}
\newtheorem{lemma}[theorem]{Lemma}
\theoremstyle{definition}
\newtheorem{definition}[theorem]{Definition}
\newtheorem{observation}[theorem]{Observation}
\newtheorem{remark}[theorem]{Remark}
\newtheorem{example}[theorem]{Example}

\newcommand{\C}{\mathcal{C}}
\newcommand{\id}{\mathrm{id}}
\newcommand{\Z}{\mathbb{Z}}
\newcommand{\Aut}{\mathrm{Aut}}
\newcommand{\Inn}{\mathrm{Inn}}
\newcommand{\Out}{\mathrm{Out}}
\newcommand{\Hom}{\mathrm{Hom}}
\newcommand{\Ab}{\bar{\mathcal{A}}}

\title{The waist class is quadratic, not additive,\\
under band composition on $T^2$}
\author{MacBeth\\ \small Kodamai}
\date{2026-10-14}

\begin{document}
\maketitle

\begin{center}
\fbox{\begin{minipage}{0.93\textwidth}
\small
\textbf{Author:} MacBeth (Kodamai).\quad\textbf{Date:} 2026-10-14.\\
\textbf{Prepared for} Rick (referee) and Neil Ghani (PI); \textbf{CC} Robin Langer.\\
\textbf{Content commit:} \texttt{3954636} (supersedes \texttt{0edf2e6}).\\
\textbf{Provenance.} Registry node \texttt{delta-additive-in-band} \textbf{[proved]} in
registry \texttt{supply-chain-inventory-gerbe-h2}; child of
\texttt{coupled-T2-general-delta-lemma} \textbf{[proved]} and
\texttt{gerbe-two-ends-correlation} \textbf{[proved]}.
Proof file \texttt{proofs/2026-10-13-delta-additive-in-band.md};
verification script \texttt{scratch/2026-10-13-delta-additivity.py} (couplers
$D_4\times C_2$, $Q_8\times C_2$; $0$ violations on every clause below).
This note typesets a complete, registry-graded result; it introduces no new
claim beyond the cited proof file.
\end{minipage}}
\end{center}

\begin{abstract}
Fix base $C=B(\Z^2)$ (nerve $T^2$) and a structure group $K$ of nilpotency class
$\le2$, so that $M:=\Inn K$ is abelian. The general $\delta$-commutator lemma
assigns to a band $\rho\colon\Z^2\to\Out K$ --- a commuting pair of outer classes --- the
connecting class $\delta(\rho)=[[A,B]]\in H^2(T^2;M)$ for any lifts $A,B\in\Aut K$.
We determine how $\delta$ behaves under band composition. The answer is: \emph{not
additively}. A master cross-term formula (Theorem~\ref{thm:master}), an on-the-nose
identity in $\Inn K$, computes the composite commutator in full. On the trivial-action
locus $\Out^0$, where the additivity question is well-typed, this collapses to a
quadratic law $\delta(\rho_1\rho_2)=\delta(\rho_1)+\delta(\rho_2)+\chi$ with symmetric
bi-additive polar form $\chi$ (Theorem~\ref{thm:quadratic}), whence $\delta(\rho^n)=n^2\delta(\rho)$.
A dichotomy (Theorem~\ref{thm:dichotomy}) shows that \emph{for an abelian subgroup
$\Ab\subseteq\Out^0$} $\delta$ is additive on $\Ab$ iff it is identically zero there:
\emph{on the trivial-action locus there is no abelian subgroup on which $\delta$ is a
nonzero homomorphism}, so a coupled torus in $\Out^0$ is always a witness to
non-additivity. (The scope $\Ab\subseteq\Out^0$ is essential: off $\Out^0$ this can
fail --- there are abelian subgroups, $1$ for $D_4\times C_2$ and $3$ for
$Q_8\times C_2$, each of order $8$ meeting $\Out^0$ in order $4$, that contain a coupled
torus $(a,b)\subseteq\Out^0$ with $\delta(a,b)\neq0$ yet on which the induced
$\bar\delta$ is additive.) On $\Out^0$, coupling, non-splitting of the lift extension,
and non-additivity of the waist class are one and the same phenomenon; this is the
classical commutator form of the central extension $\Aut^0K$ --- the
$\Hom(\Lambda^2\Ab,M)$ component of $H^2$. All clauses are verified computationally with
$0$ violations on the two $T^2$ couplers.
\end{abstract}

\section{The question and the answer}\label{sec:question}

Fix base $C=B(\Z^2)$ (nerve $T^2$) and a structure group $K$ of nilpotency class
$\le2$, so $M:=\Inn K$ is abelian. By the general $\delta$-commutator lemma
(\texttt{coupled-T2-general-delta-lemma}, proved 2026-10-12), a band
$\rho\colon\Z^2\to\Out K$ --- i.e.\ a commuting pair $\bar A:=\rho(s)$,
$\bar B:=\rho(t)\in\Out K$ --- has connecting class
\begin{equation}\label{eq:delta-def}
\delta(\rho)=[\,[A,B]\,]\in H^2(T^2;M)=M_{\Z^2}=M\big/\big((\sigma-1)M+(\tau-1)M\big),
\end{equation}
where $A,B\in\Aut K$ are \emph{any} lifts of $\bar A,\bar B$ and
$\sigma=\bar A\cdot(-)$, $\tau=\bar B\cdot(-)$ are the induced actions of
$\bar A,\bar B$ on $M$ by conjugation.

Let $\rho_1,\rho_2$ be two bands landing in a \emph{common abelian subgroup}
$\Ab\le\Out K$ (the concrete meaning of ``a common system of commuting lifts'': their
pointwise product $\rho_1\rho_2=(\bar A_1\bar A_2,\ \bar B_1\bar B_2)$ is again a band
into $\Ab$, because $\Ab$ is abelian). \textbf{Is
$\delta(\rho_1\rho_2)=\delta(\rho_1)+\delta(\rho_2)$?}

\paragraph{Answer: (B) --- No. $\delta$ is not a homomorphism.}
The exact defect is governed by the master cross-term formula
(Theorem~\ref{thm:master}), an on-the-nose identity in $\Inn K$:
\[
[A_1A_2,\,B_1B_2]\;=\;\underbrace{[A_1,B_1]}_{=\,\delta(\rho_1)\text{ rep}}
\;+\;(\sigma_1\tau_1)\!\cdot\![A_2,B_2]\;+\;\sigma_1\!\cdot\![A_2,B_1]\;+\;\tau_1\!\cdot\![A_1,B_2].
\]
On the \emph{trivial-action locus}
$\Out^0=\ker(\Out K\to\Aut M)$ --- where the question is even well-typed, all three
classes living in the \emph{same} group $M$ --- this collapses
(Theorem~\ref{thm:quadratic}) to
\[
\boxed{\;\delta(\rho_1\rho_2)=\delta(\rho_1)+\delta(\rho_2)+\chi(\rho_1,\rho_2),\qquad
\chi(\rho_1,\rho_2)=[A_1,B_2]+[A_2,B_1]\in M\;}
\]
with $\chi$ a \emph{symmetric bi-additive} form --- the polarization of the quadratic
map $\delta$, which obeys $\delta(\rho^n)=n^2\,\delta(\rho)$. And then the
\emph{dichotomy} (Theorem~\ref{thm:dichotomy}): on any abelian $\Ab\subseteq\Out^0$,
\[
\delta\text{ is additive on bands into }\Ab\iff\delta\equiv0\text{ there}\iff
\pi^{-1}(\Ab)\le\Aut K\text{ is abelian}.
\]
\textbf{On the trivial-action locus $\Out^0$ there is no abelian subgroup on which
$\delta$ is a nonzero homomorphism.} A coupled torus in $\Out^0$ is \emph{always} a
witness to non-additivity. (This is sharp to the locus: off $\Out^0$ there exist abelian
$\Ab$ --- $1$ for $D_4\times C_2$, $3$ for $Q_8\times C_2$, each of order $8$ with
$|\Ab\cap\Out^0|=4$ --- that contain a coupled torus $(a,b)\subseteq\Out^0$,
$\delta(a,b)\neq0$, yet on which the induced $\bar\delta$ is additive, the elements
outside $\Out^0$ killing $\delta(a,b)$ in coinvariants.)

So the butterfly \S7 open edge ``composite weld obstruction $=$ composite of waist
classes?'' resolves in the negative, in the sharpest possible form: the waist class
$\delta$ is a \emph{quadratic refinement}, not a homomorphism, on the composition
monoid of bands; its polar form $\chi$ is the obstruction to additivity, and $\chi$
vanishes exactly when $\delta$ itself does.

\section{The master cross-term formula}\label{sec:master}

Throughout, $M=\Inn K\trianglelefteq\Aut K$ is abelian (class $\le2$), written
additively, and for $X\in\Aut K$, $m\in M$ we write $X\cdot m:=XmX^{-1}\in M$ ($M$
normal); this factors through $\pi(X)\in\Out K$. For a band $\rho_i$ fix lifts
$A_i,B_i\in\Aut K$ of $\rho_i(s),\rho_i(t)$ and set $\sigma_i:=\rho_i(s)\cdot(-)$,
$\tau_i:=\rho_i(t)\cdot(-)\in\Aut(M)$.

\begin{theorem}[master cross-term formula]\label{thm:master}
Let $K$ be finite of class $\le2$ and let $\rho_1,\rho_2$ be bands with all four outer
classes $\bar A_1,\bar B_1,\bar A_2,\bar B_2$ in a common abelian subgroup of $\Out K$.
Then every pairwise commutator $[A_i,B_j]$ lies in $M$, and, additively in $M$,
\begin{equation}\label{eq:star}
[A_1A_2,\,B_1B_2]=[A_1,B_1]+(\sigma_1\tau_1)\!\cdot\![A_2,B_2]+\sigma_1\!\cdot\![A_2,B_1]+\tau_1\!\cdot\![A_1,B_2].\tag{$\star$}
\end{equation}
Since $A_1A_2,\,B_1B_2$ are lifts of $(\rho_1\rho_2)(s),(\rho_1\rho_2)(t)$, the left
side represents $\delta(\rho_1\rho_2)\in M_{\Z^2}$ (coinvariants for the product action
$\sigma_1\sigma_2,\tau_1\tau_2$).
\end{theorem}

\begin{proof}
Membership: $\pi[A_i,B_j]=[\bar A_i,\bar B_j]=1$ because $\bar A_i,\bar B_j$ commute
(common abelian subgroup), so $[A_i,B_j]\in\ker\pi=M$. Use the universal commutator
identities
\[
[xy,z]=x[y,z]x^{-1}\cdot[x,z],\qquad [x,yz]=[x,y]\cdot y[x,z]y^{-1}.
\]
Apply the first with $x=A_1,y=A_2,z=B_1B_2$:
\[
[A_1A_2,B_1B_2]=A_1[A_2,B_1B_2]A_1^{-1}\cdot[A_1,B_1B_2].
\]
Apply the second to each factor:
\[
[A_2,B_1B_2]=[A_2,B_1]\cdot B_1[A_2,B_2]B_1^{-1},\qquad
[A_1,B_1B_2]=[A_1,B_1]\cdot B_1[A_1,B_2]B_1^{-1}.
\]
Substituting and conjugating the first bracket by $A_1$:
\[
[A_1A_2,B_1B_2]=\big(A_1[A_2,B_1]A_1^{-1}\big)\cdot\big(A_1B_1[A_2,B_2]B_1^{-1}A_1^{-1}\big)\cdot[A_1,B_1]\cdot\big(B_1[A_1,B_2]B_1^{-1}\big).
\]
All four factors lie in the abelian group $M$, so we may reorder and pass to additive
notation. The conjugating prefixes are $A_1$ (giving $\sigma_1$), $A_1B_1$ (giving
$\sigma_1\tau_1$), nothing (giving $1$), and $B_1$ (giving $\tau_1$). This is exactly
$(\star)$.
\end{proof}

\paragraph{Verification.}
Script \texttt{scratch/2026-10-13-delta-additivity.py}, clause \texttt{[GEN]}:
$(\star)$ holds \emph{on the nose in $\Inn K$} over \emph{all} quadruples of bands
inside a common abelian subgroup --- $8113$ checks for $D_4\times C_2$ and $26037$ for
$Q_8\times C_2$, \textbf{$0$ violations}. (An earlier run comparing against the
representative lift of the \emph{product coset} showed spurious ``violations''; those
are precisely coboundaries $(\sigma-1)M+(\tau-1)M$ --- the product-coset representative
differs from $A_1A_2$ by an inner automorphism --- and vanish in $M_{\Z^2}$. Using the
on-the-nose product lift $A_1A_2$ removes them, confirming $(\star)$ is an identity in
$M$, not merely in the quotient.)

\paragraph{Reading of $(\star)$.}
Beyond the two ``diagonal'' terms $\delta(\rho_1),\delta(\rho_2)$ there are: a
\emph{twist} $(\sigma_1\tau_1)$ acting on $\delta(\rho_2)$'s representative, and the two
\emph{cross terms} $\sigma_1[A_2,B_1]+\tau_1[A_1,B_2]$. Both the twist and the asymmetry
($\sigma_1,\tau_1$ but not $\sigma_2,\tau_2$) are artifacts of writing the product lift
as $A_1A_2$ rather than $A_2A_1$; they disappear precisely when the action is trivial.

\section{The trivial-action locus: \texorpdfstring{$\delta$}{delta} is a quadratic map}\label{sec:quadratic}

\begin{definition}\label{def:out0}
$\Out^0:=\ker\!\big(\Out K\to\Aut(M)\big)$ is the subgroup of outer classes acting
trivially on $M=\Inn K$. A band $\rho$ is \emph{trivial-action} iff
$\bar A,\bar B\in\Out^0$, i.e.\ $\sigma=\tau=\id$.
\end{definition}

\begin{remark}[terminology]
The PROVE brief's ``trivial centre-action'' is here made precise as the action relevant
to $\delta$: trivial action of $\Out K$ on the coefficient module $M=\Inn K$ --- the
exact condition under which $\delta$'s receptacle $M_{\Z^2}$ collapses to $M$ and the
additivity question becomes well-typed. For class-$\le2$ $K$ this is also the condition
making the restricted extension $(\ast)|_{\Out^0}$ central, below. We do \emph{not}
silently conflate it with the distinct condition ``$\Out K$ acts trivially on $Z(K)$'';
the relevant module for the waist class is $\Inn K$, not $Z(K)$.
\end{remark}

Over $\Out^0$ the restricted extension
$1\to M\to\Aut^0K\to\Out^0\to1$ (with $\Aut^0K:=\pi^{-1}(\Out^0)$) is a \emph{central}
extension: $M$ is central in $\Aut^0K$ (verified: clause
\texttt{restricted (*) central: True} for both couplers). For such a central extension
the \emph{commutator pairing}
\[
c\colon\{\text{commuting pairs in }\Out^0\}\to M,\qquad c(\bar x,\bar y):=[\tilde x,\tilde y]
\]
is well defined (lift-independent: for central $u,v\in M$,
$[\tilde xu,\tilde yv]=[\tilde x,\tilde y]$) and $\delta(\rho)=c(\rho(s),\rho(t))$ for
trivial-action bands, with $M_{\Z^2}=M$ (no coboundaries).

\begin{lemma}[$c$ is alternating and bi-multiplicative]\label{lem:c}
On commuting pairs in $\Out^0$: $c(\bar x,\bar x)=0$; $c(\bar y,\bar x)=-c(\bar x,\bar y)$;
and whenever all pairs in sight commute,
\[
c(\bar x\bar y,\bar z)=c(\bar x,\bar z)+c(\bar y,\bar z),\qquad
c(\bar x,\bar y\bar z)=c(\bar x,\bar y)+c(\bar x,\bar z).
\]
\end{lemma}

\begin{proof}
Alternating/antisymmetry are immediate from $[\tilde x,\tilde x]=1$,
$[\tilde y,\tilde x]=[\tilde x,\tilde y]^{-1}$. For bi-multiplicativity,
\[
[\tilde x\tilde y,\tilde z]=\tilde x[\tilde y,\tilde z]\tilde x^{-1}\,[\tilde x,\tilde z]
=[\tilde y,\tilde z][\tilde x,\tilde z],
\]
using that $[\tilde y,\tilde z]\in M$ is \emph{central} (this is exactly the
trivial-action hypothesis).
\end{proof}

\noindent (Clause \texttt{(2) bi-mult on Out\^{}0}: $0$ violations; and \texttt{(2)}
over \emph{all} of $\Out$ \textbf{fails} --- $112/1408$ and $962/4032$ violations ---
confirming centrality is essential.)

\begin{theorem}[quadratic law + cross term]\label{thm:quadratic}
For trivial-action bands $\rho_1,\rho_2$ into a common abelian $\Ab\subseteq\Out^0$,
equality holds in $M$:
\[
\delta(\rho_1\rho_2)=\delta(\rho_1)+\delta(\rho_2)+\chi(\rho_1,\rho_2),\qquad
\chi(\rho_1,\rho_2)=c(\bar A_1,\bar B_2)+c(\bar A_2,\bar B_1)=[A_1,B_2]+[A_2,B_1].
\]
$\chi$ is symmetric ($\chi(\rho_1,\rho_2)=\chi(\rho_2,\rho_1)$) and bi-additive;
$\delta$ is a \emph{quadratic map}
$\Hom(\Z^2,\Ab)\cong\Ab\times\Ab\to M$ with polar form $\chi$, and in particular
$\delta(\rho^n)=n^2\,\delta(\rho)$ and $\delta(\rho^{-1})=\delta(\rho)$.
\end{theorem}

\begin{proof}
Either specialize $(\star)$ to $\sigma_1=\tau_1=\id$, or compute directly with
Lemma~\ref{lem:c}:
\[
\delta(\rho_1\rho_2)=c(\bar A_1\bar A_2,\bar B_1\bar B_2)
=c(\bar A_1,\bar B_1)+c(\bar A_1,\bar B_2)+c(\bar A_2,\bar B_1)+c(\bar A_2,\bar B_2)
=\delta(\rho_1)+\delta(\rho_2)+\chi.
\]
Symmetry of $\chi$ is visible; bi-additivity is bi-multiplicativity of $c$.
Quadraticity: $\delta(\rho^n)=c(\bar A^n,\bar B^n)=n^2c(\bar A,\bar B)=n^2\delta(\rho)$ by
bi-multiplicativity; $\delta(\rho^{-1})=c(\bar A^{-1},\bar B^{-1})=(-1)(-1)c(\bar A,\bar B)=\delta(\rho)$.
\end{proof}

\paragraph{Verification.}
Clause \texttt{(3) cross-term on Out\^{}0} ($1905$ checks each) and \texttt{(4)
quadratic} ($320$ checks each, $n\le6$): \textbf{$0$ violations} on both couplers.

\section{The additivity dichotomy}\label{sec:dichotomy}

\begin{theorem}[dichotomy]\label{thm:dichotomy}
Let $\Ab\subseteq\Out^0$ be abelian. The following are equivalent:
\begin{enumerate}
\item $\delta$ is a group homomorphism on $\Hom(\Z^2,\Ab)$ (additivity);
\item $\delta\equiv0$ on all bands into $\Ab$;
\item $c|_{\Ab\times\Ab}\equiv0$;
\item $\pi^{-1}(\Ab)\le\Aut K$ is abelian (i.e.\ $\Ab$ has a system of commuting lifts).
\end{enumerate}
In particular $\delta$ is \emph{never a nonzero homomorphism} on any abelian subgroup.
\end{theorem}

\begin{proof}
$(2)\Rightarrow(1)$ trivial. $(1)\Rightarrow(3)$: additivity forces $\chi\equiv0$; take
$\rho_1=(\bar x,1),\rho_2=(1,\bar y)$, so $\delta(\rho_1)=c(\bar x,1)=0$,
$\delta(\rho_2)=c(1,\bar y)=0$, while $\rho_1\rho_2=(\bar x,\bar y)$ gives
$\delta(\rho_1\rho_2)=c(\bar x,\bar y)$; additivity makes this $0$ for all
$\bar x,\bar y\in\Ab$. $(3)\Rightarrow(2)$: $\delta(\rho)=c(\bar A,\bar B)=0$.
$(3)\Leftrightarrow(4)$: $\pi^{-1}(\Ab)$ is a central extension of the abelian $\Ab$ by
$M$, hence abelian iff its commutator pairing $c|_{\Ab}$ vanishes.
\end{proof}

\paragraph{Verification.}
Clause \texttt{(5) dichotomy on Out\^{}0}: among the $6$ coupled abelian subgroups of
$\Out^0$ (for each coupler), the count \texttt{nonzero \& ADDITIVE} is \textbf{$0$};
every coupled subgroup is non-additive. Exhaustive over $\le2$-generated abelian
subgroups (which suffices, bands being rank $\le2$).

\paragraph{Reconciliation with the coupling criterion.}
$(4)$ is the negation of the $\delta$-lemma's coupling criterion ``$\exists$ commuting
outer pair with no commuting lift.'' So Theorem~\ref{thm:dichotomy} says, \emph{for
abelian $\Ab\subseteq\Out^0$}:
\[
\Ab\text{ couples (some band has }\delta\neq0)\iff\pi^{-1}(\Ab)\text{ nonabelian}
\iff\delta\text{ is non-additive on }\Ab.
\]
On the trivial-action locus, coupling, non-splitting of $\pi^{-1}(\Ab)$, and
non-additivity of the waist class are \emph{one and the same phenomenon} --- this is the
classical commutator form of a central extension, the $\Hom(\Lambda^2\Ab,M)$ component
of $H^2(\Ab;M)$, and we claim no novelty for it here. Off $\Out^0$ the chain breaks at
$(4)$: $\pi^{-1}(\Ab)$ need not be central over $\Ab$, and (Theorem~\ref{thm:receptacle})
an abelian $\Ab\not\subseteq\Out^0$ can contain a coupled torus while $\bar\delta$ is
additive on it ($1$ such $\Ab$ for $D_4\times C_2$, $3$ for $Q_8\times C_2$).

\section{Scope, and the genuinely twisted regime}\label{sec:scope}

The clean additivity package (Theorems~\ref{thm:quadratic}--\ref{thm:dichotomy}) lives
on the \emph{trivial-action locus} $\Out^0$, and this is not a cosmetic restriction but
the precise place the question is \emph{well-typed}: for bands with differing actions,
$\delta(\rho_1),\delta(\rho_2),\delta(\rho_1\rho_2)$ lie in \emph{different} coinvariant
quotients $M_{(\sigma,\tau)}$, so ``$\delta(\rho_1\rho_2)=\delta(\rho_1)+\delta(\rho_2)$''
is not even a statement about one group until a common receptacle is chosen.

\paragraph{Coupling is \emph{not} confined to $\Out^0$.}
The census (\texttt{coupling vs action}) finds coupled bands ($\delta\neq0$ in the
quotient) with \emph{nontrivial} action: $24$ of $60$ for $D_4\times C_2$, $72$ of $108$
for $Q_8\times C_2$. For these, the governing statement is the master formula $(\star)$
alone --- a correct on-the-nose identity in $\Inn K$. Its descent to a common
coinvariant quotient does \emph{not}, however, retain the twist: in the full
coinvariant module the factor $(\sigma_1\tau_1-1)\!\cdot\!\delta(\rho_2)$ lies in the
augmentation submodule $\sum_x(x-1)M$, so the twist \emph{dies} in that quotient (it is
computed nonzero in $0/10304$ and $0/36420$ checks). The honest common receptacle is
therefore the subgroup-coinvariant module itself:
\[
M_{\Ab}\;:=\;M\Big/\textstyle\sum_{x\in\Ab}(x-1)M,
\]
the coinvariants for the whole abelian subgroup $\Ab=\langle\bar A_1,\bar B_1,\bar A_2,
\bar B_2\rangle$ generated by the four outer classes.

\begin{theorem}[common receptacle off $\Out^0$]\label{thm:receptacle}
Let $\Ab\le\Out K$ be abelian (not necessarily $\subseteq\Out^0$). In $M_{\Ab}$ the
induced commutator pairing $\bar c$ is well defined and bi-additive, and
Theorem~\ref{thm:quadratic} together with the equivalences
$(1)\Leftrightarrow(2)\Leftrightarrow(3)$ of Theorem~\ref{thm:dichotomy} hold
\emph{verbatim} with $\bar c,\bar\delta$ in place of $c,\delta$. The remaining
equivalence ``$(1)\text{--}(3)\Leftrightarrow(4)$'' $[\,\pi^{-1}(\Ab)$ abelian$\,]$ is
the one clause that \emph{fails} off $\Out^0$.
\end{theorem}

\noindent Verified: $0$ violations on \emph{all} $14$ abelian subgroups of
$\Out(D_4\times C_2)$ and all $50$ of $\Out(Q_8\times C_2)$ that are not
$\subseteq\Out^0$. This closes --- rather than defers --- the common-receptacle edge: the
only genuine loss off $\Out^0$ is clause $(4)$, since the lift extension
$\pi^{-1}(\Ab)$ need no longer be a central extension of $\Ab$.
\begin{itemize}
\item \textbf{No gap in the posed question.} The question fixes trivial centre-action;
there the answer is complete and proved: (B), with $\chi=[A_1,B_2]+[A_2,B_1]$, quadratic
law, dichotomy.
\end{itemize}

\section{Verification summary (hostile-referee pass)}\label{sec:verif}

All on \texttt{scratch/2026-10-13-delta-additivity.py}, couplers $D_4\times C_2$
($|\!\Out|=16$), $Q_8\times C_2$ ($|\!\Out|=48$):

\begin{center}
\renewcommand{\arraystretch}{1.3}
\begin{tabular}{>{\raggedright}p{0.17\textwidth}>{\raggedright}p{0.40\textwidth}>{\raggedright}p{0.21\textwidth}>{\raggedright\arraybackslash}p{0.14\textwidth}}
\hline
\textbf{clause} & \textbf{statement} & \textbf{checks ($D_4{\times}C_2$ / $Q_8{\times}C_2$)} & \textbf{violations}\\
\hline
\texttt{[GEN]} & master formula $(\star)$ on-the-nose in $M$ & 8113 / 26037 & \textbf{0 / 0}\\
central & $(\ast)$ restricted to $\Out^0$ is central & --- & True / True\\
\texttt{(2)} & $c$ bi-multiplicative on $\Out^0$ & 512 / 512 & \textbf{0 / 0}\\
\texttt{(2)} over all $\Out$ & bi-multiplicativity \textbf{fails} off $\Out^0$ & 1408 / 4032 & 112 / 962 \emph{(expected)}\\
\texttt{(3)} & $\delta(\rho_1\rho_2)=\delta_1+\delta_2+\chi$ on $\Out^0$ & 1905 / 1905 & \textbf{0 / 0}\\
\texttt{(4)} & $\delta(\rho^n)=n^2\delta(\rho)$ on $\Out^0$ & 320 / 320 & \textbf{0 / 0}\\
\texttt{(5)} & dichotomy: nonzero \& additive & 6 coupled subs each & \textbf{0 / 0}\\
\hline
\end{tabular}
\end{center}

\begin{itemize}
\item \textbf{Theorem~\ref{thm:master} --- GREEN} (hand proof from universal commutator
identities $+$ $M$ abelian; computed confirmation $0$ violations). The only hypotheses
are class $\le2$ and common abelian subgroup.
\item \textbf{Theorems~\ref{thm:quadratic}--\ref{thm:dichotomy} --- GREEN} (two-line
proofs from the central commutator pairing; computed confirmation).
\item \textbf{Boundary cases.} Trivial band $\rho=1$:
$\chi(1,\rho)=c(1,\bar B)+c(\bar A,1)=0$, additive $\checkmark$. $\rho_1=\rho_2$:
$\chi(\rho,\rho)=2c(\bar A,\bar B)=2\delta(\rho)$, consistent with
$\delta(\rho^2)=4\delta(\rho)$ $\checkmark$.
\item \textbf{No RED. No YELLOW} inside the stated scope; the twisted edge of
\S\ref{sec:scope} is flagged and is not part of the posed question.
\end{itemize}

\section{Grant / program placement}\label{sec:grant}

This is the \emph{composition-monoid structure of the waist class}. The butterfly /
weak-morphism arc asked whether the weld obstruction composes additively; the answer
sharpens the whole gerbe lane:
\begin{itemize}
\item \textbf{Theory.} The waist class $\delta$ is a \emph{quadratic refinement} on the
band composition monoid, with explicit polar form $\chi$ $=$ the commutator pairing of
the central extension $\Aut^0K$. This is the $T^2$ counterpart of the floor $D_1/D_2$
dichotomy and of the two-ends forced correlation: ``coupling $=$ nonabelian lift $=$
non-additivity'' is a single statement.
\item \textbf{Applications.} For orchestration / ZS welds, composing two welds does
\emph{not} add their obstruction classes; the correction $\chi$ is an economically real
interaction term (the ``cross-coupling'' of two independently-welded subsystems).
Additivity of the composite, $\delta(\rho_1\rho_2)=\delta(\rho_1)+\delta(\rho_2)$, holds
iff the cross-term $\chi(\rho_1,\rho_2)=c(\bar A_1,\bar B_2)+c(\bar A_2,\bar B_1)$
vanishes on the \emph{pair} --- \emph{not} iff each subsystem is individually uncoupled.
Indeed $\rho_1=(\bar x,1)$, $\rho_2=(1,\bar y)$ have $\delta(\rho_1)=\delta(\rho_2)=0$
yet $\delta(\rho_1\rho_2)=c(\bar x,\bar y)$, nonzero whenever $\bar x,\bar y$ couple: two
individually-uncoupled welds can interact.
\item \textbf{Formalisation.} Theorem~\ref{thm:master}'s identity $(\star)$ and
Lemma~\ref{lem:c} are finite, decidable, and a natural Lean target over the existing
$\Inn/\Out$ engine.
\end{itemize}

\paragraph{One-line summary.}
On $T^2$ the band-to-waist connecting map is a \emph{quadratic map}, not a homomorphism:
in general
$[A_1A_2,B_1B_2]=[A_1,B_1]+(\sigma_1\tau_1)[A_2,B_2]+\sigma_1[A_2,B_1]+\tau_1[A_1,B_2]$,
which on the trivial-action locus reads
$\delta(\rho_1\rho_2)=\delta(\rho_1)+\delta(\rho_2)+\big([A_1,B_2]+[A_2,B_1]\big)$ with
symmetric bi-additive polar form $\chi$; and \emph{for an abelian subgroup
$\Ab\subseteq\Out^0$} $\delta$ is additive on $\Ab$ \textbf{iff} it is identically zero
there --- a coupled torus in $\Out^0$ is always a witness to non-additivity. (Off
$\Out^0$ the induced $\bar\delta$ can be additive on an abelian subgroup that still
contains a coupled torus: $1$ such subgroup for $D_4\times C_2$, $3$ for
$Q_8\times C_2$.)

\end{document}
